\par\noindent\textbf{(1)} 用 \(\sigma_-=|0\rangle\langle1|\)、\(\sigma_+=|1\rangle\langle0|\)，主方程 \[\dot\rho=-i[H,\rho]/\hbar+\Gamma_\downarrow\mathcal D[\sigma_-]\rho+\Gamma_\uparrow\mathcal D[\sigma_+]\rho\]，其中 \(\mathcal D[L]\rho=L\rho L^\dagger-\{L^\dagger L,\rho\}/2\)。
\par\noindent\textbf{(2)} 激发态布居 \(p_1=\rho_{11}\) 满足 \(\dot p_1=-\Gamma_\downarrow p_1+\Gamma_\uparrow(1-p_1)\)，解为 \(p_1(t)=p_1^{\rm eq}+[p_1(0)-p_1^{\rm eq}]e^{-(\Gamma_\uparrow+\Gamma_\downarrow)t}\)，\(p_1^{\rm eq}=\Gamma_\uparrow/(\Gamma_\uparrow+\Gamma_\downarrow)\)。
\par\noindent\textbf{(3)} 利用热详细平衡 \(\Gamma_\uparrow/\Gamma_\downarrow=e^{-\beta\hbar\omega_0}\)，得到 \(p_1^{\rm eq}=1/(1+e^{\beta\hbar\omega_0})\)，恰与两能级Gibbs态 \(\rho_{\rm G}=e^{-\beta H}/\operatorname{tr}e^{-\beta H}\) 一致。若驱动浴不满足详细平衡，则固定点不一定为热态。
\par\noindent\textbf{(4)} 非对角元 \(\dot\rho_{01}=[i\omega_0-(\Gamma_\uparrow+\Gamma_\downarrow)/2]\rho_{01}\)，故 \(T_2=2/(\Gamma_\uparrow+\Gamma_\downarrow)=2T_1\)，其中 \(T_1=1/(\Gamma_\uparrow+\Gamma_\downarrow)\)。若另有纯退相干 \(\gamma_\varphi\)，则 \(1/T_2=1/(2T_1)+\gamma_\varphi\)。
